\documentclass[11pt]{ctexart}
%=====================================================================
% theme-violet.tex —— 讲义精讲 · 紫花主题
% 用法：\documentclass[11pt]{ctexart}
%       \input{assets/theme-violet}
%       \begin{document} ... \end{document}
% 只用 ctexart 自带能力 + xcolor/enumitem/fancyhdr/pifont/amsmath，
% 不依赖 tikz/tcolorbox，保证 MiKTeX 可直接编译。
%=====================================================================

\usepackage[a4paper,left=2.0cm,right=2.0cm,top=2.1cm,bottom=2.3cm,%
            headsep=12pt,footskip=22pt]{geometry}
\usepackage{amsmath,amssymb}
\usepackage{needspace}
\usepackage{xcolor}
\usepackage{array}
\usepackage{enumitem}
\usepackage{pifont}
\usepackage{fancyhdr}
\usepackage[hidelinks]{hyperref}

%---------------------------------------------------------------- 配色
\definecolor{violet}{RGB}{107,78,143}      % 主色：紫
\definecolor{violetdeep}{RGB}{63,44,92}    % 深紫：标题/强调
\definecolor{petal}{RGB}{183,159,208}      % 花瓣：装饰线、花饰
\definecolor{bandviolet}{RGB}{247,243,251} % 浅紫底：档案/提示框
\definecolor{bandcream}{RGB}{252,249,242}  % 米白底：易错/注意框
\definecolor{rulepurple}{RGB}{217,207,230} % 分隔线
\definecolor{inkgray}{RGB}{92,96,110}      % 次要文字

%---------------------------------------------------------------- 字体
% 有 Noto CJK 用 Noto，没有就退回 ctex 默认（Windows 下为宋体/黑体）
\IfFontExistsTF{Noto Serif CJK SC}{\setCJKmainfont{Noto Serif CJK SC}}{}
\IfFontExistsTF{Noto Sans CJK SC}{\setCJKsansfont{Noto Sans CJK SC}}{}

%---------------------------------------------------------------- 版式
\setlist{nosep,leftmargin=1.9em}
\linespread{1.28}
\setlength{\parskip}{3.5pt}
\setlength{\parindent}{0pt}
\setlength{\tabcolsep}{5pt}
\renewcommand{\arraystretch}{1.28}

%---------------------------------------------------------------- 花饰
% \ding{95} = 五瓣小花，\ding{96} = 星形花；都来自 pifont，必定存在
\newcommand{\orn}{\textcolor{petal}{\ding{95}}}          % 行内小标
\newcommand{\ornstar}{\textcolor{petal}{\ding{96}}}      % 星形花
\newcommand{\flowerrule}{%                                % 居中花饰分隔线
  \par\vspace{7pt}%
  \noindent\makebox[\linewidth]{\color{petal}%
    \rule{0.20\linewidth}{0.6pt}\hspace{8pt}{\small\ding{95}}\hspace{4pt}%
    \raisebox{1.1pt}{\tiny\ding{95}}\hspace{4pt}{\small\ding{95}}\hspace{8pt}%
    \rule{0.20\linewidth}{0.6pt}}%
  \par\vspace{7pt}}

%---------------------------------------------------------------- 标题与节头
\newcommand{\doctitle}[3]{%   #1 主标题  #2 副标题  #3 元信息（可用 \\ 分行）
  \begin{center}
    {\color{petal}\small \orn\enspace #2\enspace \ornstar}\\[7pt]
    {\LARGE\bfseries\color{violetdeep} #1}\\[10pt]
    {\color{petal}\rule{0.55\linewidth}{1pt}}\\[4pt]
    {\small\color{inkgray} #3}
  \end{center}\vspace{2pt}}

\newcommand{\hdr}[2]{%   #1 章节号（如 §0）  #2 标题
  \par\vspace{12pt}%
  \noindent\colorbox{violetdeep}{\color{white}\bfseries\small\;#1\;}\hspace{7pt}%
  {\color{violetdeep}\large\bfseries #2}\par\nobreak
  \vspace{3pt}{\color{petal}\hrule height 1.2pt}\vspace{6pt}}

\newcommand{\subhdr}[1]{%
  \par\vspace{6pt}%
  \noindent{\color{violet}\bfseries #1}\par\vspace{1pt}}

%---------------------------------------------------------------- 提示框
\newcommand{\callout}[1]{% 浅紫底：讲义档案、读法说明
  \par\medskip\noindent
  \fcolorbox{rulepurple}{bandviolet}{\parbox{\dimexpr\linewidth-2\fboxsep-2\fboxrule}{\vspace{2pt}#1\vspace{2pt}}}%
  \par\medskip}

\newcommand{\warnbox}[1]{% 米白底：易错、边界条件
  \par\medskip\noindent
  \fcolorbox{rulepurple}{bandcream}{\parbox{\dimexpr\linewidth-2\fboxsep-2\fboxrule}{\vspace{2pt}#1\vspace{2pt}}}%
  \par\medskip}

\newcommand{\gapitem}[2]{% #1 缺口类型（紫） #2 内容
  \par\vspace{2pt}\noindent
  \hangindent=1.4em\hangafter=1
  {\color{petal}\small\ding{96}}\enspace{\bfseries\color{violetdeep}#1}\enspace #2\par}

% 题干块（\qstem）：每道例题／作业的「详解」开头必须先复述题干与出处。
% 起因：2026-10-25 第五章讲义把七道作业的题面集中放在一页、详解另起一页，
% 读者翻到「第 4 题」时看不到条件，等于一道没有题干的题目。
% 规则：详解块内自带题干，读者不翻页就能读懂题在问什么。
\newcommand{\qstem}[1]{%
  \par\medskip\noindent
  \fcolorbox{violet}{bandviolet}{%
    \parbox{\dimexpr\linewidth-2\fboxsep-2\fboxrule}{%
      \vspace{2pt}%
      {\small\bfseries\color{violet}题干}\hspace{0.8em}#1%
      \vspace{2pt}}}%
  \par\medskip}

% v2: 保护题干与开头，长题可跨页，续页补上下文。
\newcommand{\qbreak}{\needspace{6\baselineskip}}
% 独立练习：ID、来源、题干、考点、答案标签（通常 ans:ID）。
\newcommand{\practice}[5]{%
  \qbreak\subhdr{练习 #1}\qstem{（#2）#3}%
  \textbf{考点}：#4\par
  详解见附录第~\pageref{#5}~页。\par\medskip}
% 附录开头：ID、来源、条件摘要；自动产生 ans:ID 标签。
\newcommand{\answerstart}[3]{%
  \qbreak\subhdr{#1\quad 详解}\label{ans:#1}\qstem{（#2）#3}}
\newcommand{\answercontinue}[2]{%
  \needspace{5\baselineskip}\subhdr{#1\quad 续解}%
  {\small\color{inkgray}本页接续：#2}\par\medskip}
% 延迟提供默认值，让旧讲义可在导言区保留自己的 \newcommand{\step}。
\AtBeginDocument{\providecommand{\step}[1]{\needspace{3\baselineskip}\par\medskip\textbf{#1}\par}}

% 前提旁注：只放在「需要理解的那一步」后面，给陈述 + 本题代入，不重证。
\newcommand{\aside}[1]{%
  \par\vspace{1pt}\noindent
  \hangindent=1.6em\hangafter=1
  {\small\color{inkgray}%
    \hspace*{0.6em}{\color{petal}$\triangleright$}\hspace{2pt}#1}%
  \par\vspace{2pt}}

%---------------------------------------------------------------- 表格
% 本机 MiKTeX 的 colortbl 与 array 版本不匹配（\rowcolor 会报错），
% 故表格统一用「上粗紫线 + 表头 + 中细线」的三线式，不依赖 colortbl。
\newcommand{\tabtop}{\par\vspace{3pt}{\color{violet}\hrule height 1.2pt}\vspace{1pt}}
\newcommand{\tabmid}{\par\vspace{1pt}{\color{petal}\hrule height 0.5pt}\vspace{1pt}}
\newcommand{\tabbot}{\par\vspace{3pt}{\color{rulepurple}\hrule height 0.5pt}\vspace{8pt}}
\newcommand{\thead}[1]{{\bfseries\color{violetdeep}#1}}

%---------------------------------------------------------------- 页眉页脚
\pagestyle{fancy}
\fancyhf{}
\renewcommand{\headrulewidth}{0.4pt}
\renewcommand{\headrule}{\hbox to\headwidth{\color{rulepurple}\leaders\hrule height \headrulewidth\hfill}}
\fancyhead[L]{\small\color{inkgray}\songti \doctag}
\fancyhead[R]{\small\color{inkgray}\songti \coursetag}
\fancyfoot[C]{\small\color{inkgray}\orn\enspace{\color{violet}\thepage}\enspace\orn}

% 每篇讲义在导言区设这两个值
\newcommand{\doctag}{}
\newcommand{\coursetag}{}
\newcommand{\settags}[2]{\renewcommand{\doctag}{#1}\renewcommand{\coursetag}{#2}}

%---------------------------------------------------------------- 数学记号
\newcommand{\vect}[1]{\boldsymbol{#1}}
\newcommand{\dd}{\,\mathrm{d}}
\newcommand{\ddx}{\frac{\mathrm{d}}{\mathrm{d}x}}
\addtolength{\jot}{2pt}



\geometry{headheight=22pt}

% ---- 本讲义专用记号 ----
\newcommand{\pmat}[1]{\begin{pmatrix}#1\end{pmatrix}}
\newcommand{\R}{\mathbb{R}}
\newcommand{\ii}{\mathrm{i}}
\newcommand{\zero}{[0]}
\newcommand{\Ax}{Ax}

\settags{矩阵指数}{微分方程}
\begin{document}

\doctitle{矩阵指数与常系数线性方程组}{精讲讲义 · Matrix Exponential}%
{源资料：Differential Equations，Lecture 16--17（W.\,Y., DLUT，2026-10-13，共 44 页）\\
覆盖：p1--p44；课件全部 8 个编号例题（Example 1--8）与 p44 的 3 道作业\\
用途：可独立自学的精读讲义；讲义覆盖不等于已学或已验收}

\flowerrule

%=====================================================================
\hdr{§0}{读前须知：对象、记号与自检}
%=====================================================================

\subhdr{0.1 这一讲解决什么问题}

标量初值问题
\[
x'=ax,\qquad x(0)=x_0,\qquad a\in\R\ \text{为常数}
\]
的解是 $x=x_0e^{at}$。本讲把同一个问题搬到方程组：
\begin{equation}\label{eq:system}
\frac{\dd Y}{\dd x}=AY,\qquad Y(0)=Y_0,\qquad x\in\R,
\end{equation}
其中 $A$ 是 $n\times n$ 常数矩阵，$Y$ 是 $n$ 维列向量。自然要问：解是否就是
$Y(x)=e^{Ax}Y_0$？答案是“是”，但前提是先把 $e^{Ax}$ 定义成一个 $n\times n$ 矩阵，
并证明它确实满足方程。课件随后给出两条实际计算 $e^{Ax}$ 的路线：
\begin{itemize}
  \item 当 $A$ 有“足够多”的线性无关特征向量时，用对角化
        $e^{Ax}=Ve^{\Lambda x}V^{-1}$（课件 §2，p8--p15）；
  \item 当特征向量不够（有重特征值）时，用 Cayley--Hamilton 定理把无穷级数压成
        关于 $A$ 的次数不超过 $n-1$ 的多项式（课件 §3，p16--p38）。
\end{itemize}
最后用“基本矩阵”说明 $e^{Ax}$ 与一般线性系统的解矩阵之间的关系（p39--p43）。
本讲的结构与课件目录（p2）一致；p3、p8、p16、p39 是各部分的过渡页。

\clearpage
\subhdr{0.2 记号表}

\begin{center}
\begin{tabular}{p{0.14\linewidth}p{0.76\linewidth}}
\hline
\thead{记号} & \thead{含义（首次出现处）}\\
\hline
$E$ & $n\times n$ 单位矩阵；$A^0=E$（p4--p5）\\
$\zero$ & $n\times n$ 零矩阵（p5；课件写作 $[0]$）\\
$e^{Ax}$ & 矩阵指数，一个 $n\times n$ 矩阵，不是数（p5）\\
$V_j,\ \lambda_j$ & $A$ 的特征向量与特征值（p9）；$V=(V_1,\dots,V_n)$\\
$\Lambda$ & $\Lambda=V^{-1}AV=\operatorname{diag}(\lambda_1,\dots,\lambda_n)$（p9--p10）\\
$m_j,\ g_j$ & 特征值 $\lambda_j$ 的代数重数与几何重数（p9）\\
$\Phi(x)$ & 基本矩阵，列是基本解组（p40）\\
$\Delta(s)$ & 首一特征多项式 $\det(sE-A)$ 或 $\det(sE-\Ax)$；与 $\det(A-sE)$ 相差 $(-1)^n$（p19、p22）\\
$r(\lambda)$ & 次数 $\le n-1$ 的多项式 $\sum_{k=0}^{n-1}\alpha_k\lambda^k$（p21--p22）\\
$\alpha_k$ & $r$ 的系数；在 $e^{Ax}$ 的计算中它们是 $x$ 的函数（p18）\\
$C,\ Y_0$ & 任意常向量与初值向量（p5、p6）\\
\hline
\end{tabular}
\end{center}

\aside{记号冲突提醒：课件在 p9--p12 用 $\lambda_j$ 表示 \emph{$A$ 的特征值}；在 p22 以后又用
$\lambda_i$ 表示 \emph{$\Ax$ 的特征值}。两者相差一个因子 $x$。本讲义在 §3 沿用前一种用法；
在 §4 把 $A$ 的特征值改记作 $\mu$，把 $\Ax$ 的特征值记作 $\lambda=x\mu$，以免混淆。
课件原式中的 $\lambda_i$ 一律按“$\Ax$ 的特征值”理解。}

\subhdr{0.3 需要的前置，以及“会不会用”的判断}

讲义默认读者会矩阵乘法、行列式、特征值与特征向量、线性无关、幂级数。真正容易出错的是
“条件用在哪一步”，所以这里先给出可操作的检查方式，后文用到时会再点一次：
\begin{itemize}
  \item \textbf{矩阵乘法不可交换。} 只有当 $AB=BA$ 时，才能把 $(A+B)^n$ 按二项式展开，
        从而把 $e^{A+B}$ 拆成 $e^Ae^B$。判断方法：直接算 $AB$ 与 $BA$，看两个矩阵是否相等。
  \item \textbf{特征向量要验。} 算出候选 $V\ne0$ 后，代回 $AV=\lambda V$ 逐分量核对；
        只算 $\det(A-\lambda E)=0$ 只得到特征值，不自动得到特征向量。
  \item \textbf{线性无关要验。} 把 $V_1,\dots,V_n$ 作为列拼成 $V$，检查 $\det V\ne0$
        （或对 $V$ 做行变换看是否有 $n$ 个主元）；通过后才允许写 $V^{-1}$。
  \item \textbf{矩阵等式要整体验。} 矩阵等式不能在某个分量上“看起来对”就通过；
        要逐元素相等，或代回原式作矩阵乘法。
  \item \textbf{矩阵的“大小”（范数，可选）}：§1.3 要用一个数估计矩阵幂的增长。没学过范数时，
        只需接受那里给出的两条性质（乘法不放大太多、级数被 $e^{\|A\||x|}$ 控制）；
        后文按特征值计算并不依赖范数的具体选择。
\end{itemize}

\needspace{18\baselineskip}
\subhdr{0.4 本讲义反复使用的三条自检}

后面每个例题算完，都至少做下面三项中的两项。它们不是形式，而是最快发现符号错误的方法：
\begin{enumerate}
  \item \textbf{代 $x=0$：} $e^{A\cdot0}=E$。所以任何候选 $e^{Ax}$ 在 $x=0$ 必须等于单位矩阵。
        这一条能立刻查出系数漏乘、特征向量矩阵取错等问题。
  \item \textbf{求导后对照 $A$：} $\left.\dfrac{\dd}{\dd x}e^{Ax}\right|_{x=0}=A$。
        这一条检查指数上的系数与矩阵 $A$ 是否匹配。
  \item \textbf{逐列代回方程：} $e^{Ax}$ 的每一列都必须是 $Y'=AY$ 的解；
        等价地，检查 $\dfrac{\dd}{\dd x}e^{Ax}=Ae^{Ax}$（矩阵等式）。
\end{enumerate}

%=====================================================================
\hdr{§1}{从标量指数解到矩阵指数（p4--p6）}
%=====================================================================

\subhdr{1.1 标量情形：先确认要推广的“答案形状”}

对 $x'=ax$、$x(0)=x_0$，解是
\begin{equation}\label{eq:scalar}
x(t)=x_0e^{at}.
\end{equation}
把它写成“初值乘一个指数因子”的形状，正是后面要推广的对象。课件 p4 把这个式子排成了
$x=x_0e^{atx}$（指数上多了一个 $x$）。把 $x=x_0e^{atx}$ 代回 $x'=ax$ 会多出含 $x'$ 的项，
不满足方程；正确形式是 \eqref{eq:scalar}。这一处按勘误处理，详细核对见附录 B。

\subhdr{1.2 从数到矩阵：先排除一个错误直觉}

把 $a$ 换成矩阵 $A$、把 $x_0$ 换成向量 $Y_0$，最省事的写法是把 $e^{Ax}$ 理解成
“对 $A$ 的每个元素分别取指数”。这个直觉一般是错的。取
\[
A=\pmat{0&1\\0&0},\qquad A^2=\zero .
\]
按定义（下一小节）$e^{A}=E+A=\pmat{1&1\\0&1}$；而“逐元素取指数”会得到
$\pmat{1&e\\1&1}$，两者不同。原因是指数级数里出现的是 $A$ 的\emph{幂}，矩阵乘法会把不同位置的元素混合起来。
所以 $e^{Ax}$ 必须由级数定义，不能逐元素操作。

\subhdr{1.3 定义：矩阵指数是一个矩阵}

\callout{
\textbf{定义（课件 Theorem 1.1，p5）.} 设 $A$ 是 $n\times n$ 矩阵，定义
\[
e^{Ax}=E+Ax+\frac{A^2x^2}{2!}+\cdots+\frac{A^kx^k}{k!}+\cdots
=\sum_{k=0}^{\infty}\frac{A^kx^k}{k!},\qquad A^0=E .
\]
该级数对每个实数（或复数）$x$ 收敛，其和是一个 $n\times n$ 矩阵，称为矩阵指数。}

\textbf{为什么收敛（最小必要说明）.} 把 $n\times n$ 矩阵看成 $n^2$ 个数排成的表，用一个数衡量它
“能把向量放大多少”。下面只用到两条性质（记 $B=(b_{ij})$）：
\begin{itemize}
  \item \textbf{大小}：取 $\|B\|=\max_{1\le i\le n}\sum_{j=1}^n|b_{ij}|$（最大行和）。它满足
        $\|BC\|\le\|B\|\,\|C\|$，因此 $\|B^k\|\le\|B\|^k$。
  \item \textbf{收敛}：由上式，
        $\|A^kx^k/k!\|\le\|A\|^k|x|^k/k!$，所以
        \[
        \sum_{k=0}^{\infty}\frac{\|A^kx^k\|}{k!}
        \le\sum_{k=0}^{\infty}\frac{\|A\|^k|x|^k}{k!}
        =e^{\|A\|\,|x|}<\infty .
        \]
\end{itemize}
于是 $e^{Ax}$ 的级数在每个矩阵元素上绝对收敛。这一段只用“矩阵可以量大小、乘法不会把大小
放大得比乘积更糟”这一件事；不熟悉范数的读者可把它当作选读，后文的计算不依赖范数的具体选择。
逐项求导与重排的依据是幂级数自身的性质，见 §1.4，不是仅由“绝对收敛”三个字直接推出。

\textbf{两个大小关系要分清：} $e^{Ax}$ 是 $n\times n$ 矩阵；$e^{\lambda x}$（$\lambda$ 是数）是数。
课件特别用 NB 标出前者，正是因为初学时容易把矩阵指数当成“每个元素一个指数”。

\subhdr{1.4 求导验证：$e^{Ax}C$ 确实是解}

课件 p5 的核心计算是
\[
\left(e^{Ax}C\right)'-Ae^{Ax}C
=\frac{\dd}{\dd x}\left(\sum_{k=0}^{\infty}\frac{(\Ax)^k}{k!}\right)C
-A\sum_{k=0}^{\infty}\frac{(\Ax)^k}{k!}C .
\]
这里 $\Ax$ 表示矩阵 $A$ 乘数 $x$，$(\Ax)^k=A^kx^k$。$\sum_{k=0}^{\infty}A^kx^k/k!$ 是幂级数
（系数是矩阵），收敛半径 $\infty$；幂级数在其收敛区间内部可以逐项求导，且在任意有限区间上
一致收敛，所以下面的逐项求导合法。注意这是幂级数的性质，不能对任意函数项级数照搬。
\[
\frac{\dd}{\dd x}\sum_{k=0}^{\infty}\frac{A^kx^k}{k!}
=\sum_{k=1}^{\infty}\frac{kA^kx^{k-1}}{k!}
=\sum_{k=1}^{\infty}\frac{A^k x^{k-1}}{(k-1)!}
=A\sum_{j=0}^{\infty}\frac{A^jx^j}{j!}=Ae^{Ax}.
\]
第三步把 $kA^k=kA\cdot A^{k-1}$ 中的 $A$ 提到求和号外，并令 $j=k-1$；最后一步用回定义。
于是
\[
\left(e^{Ax}C\right)'-Ae^{Ax}C=Ae^{Ax}C-Ae^{Ax}C=\zero,
\]
即 $Y=e^{Ax}C$ 满足 $Y'=AY$。注意这里能提出 $A$，是因为 $A$ 与自己的每个幂可交换；
若把 $A$ 换成依赖 $x$ 的矩阵 $A(x)$，这个推导立刻失效（§5 会说明此时只能用基本矩阵）。

\subhdr{1.5 初值问题与“为什么这是全部解”}

若给定 $Y(x_0)=Y_0$，则
\begin{equation}\label{eq:ivp}
Y(x)=e^{A(x-x_0)}Y_0
\end{equation}
满足方程与初值：$Y(x_0)=e^{A\cdot0}Y_0=Y_0$，且
$Y'=Ae^{A(x-x_0)}Y_0=AY$。课件的展开式
$e^{A(x-x_0)}=E+A(x-x_0)+\cdots$ 只是把定义中的 $x$ 换成 $x-x_0$。

“全部解”还需要一句话：$e^{Ax}$ 可逆，逆是 $e^{-Ax}$（§2 证明）。于是对任一解 $Y$，
\[
\frac{\dd}{\dd x}\left(e^{-Ax}Y\right)
=-Ae^{-Ax}Y+e^{-Ax}Y'
=e^{-Ax}(-AY+AY)=0,
\]
所以 $e^{-Ax}Y\equiv C$ 为常向量，即 $Y=e^{Ax}C$。这同时说明：
$e^{Ax}$ 的 $n$ 个列向量构成解空间的一组基，通解正是 $Y=e^{Ax}C$。

\subhdr{1.6 教学小例：三种一眼能算的矩阵}

这三个例子说明“级数什么时候能直接求和”，也为 §3、§4 的方法提供直觉。

\textbf{例 A（对角矩阵）.} $A=\operatorname{diag}(\lambda_1,\lambda_2)$。因为
$A^k=\operatorname{diag}(\lambda_1^k,\lambda_2^k)$，逐元素求和得
\[
e^{Ax}=\pmat{e^{\lambda_1x}&0\\0&e^{\lambda_2x}}.
\]
这正是 §3 对角化公式的最简单情形。

\textbf{例 B（纯量矩阵）.} $A=rE$。则 $A^k=r^kE$，
\[
e^{Ax}=\sum_{k=0}^{\infty}\frac{r^kx^k}{k!}E=e^{rx}E .
\]
注意右边是“数 $e^{rx}$ 乘单位矩阵”，不是把 $e^{rx}$ 塞进每个元素后的任意矩阵。

\textbf{例 C（幂零矩阵）.} 取 $N=\pmat{0&1\\0&0}$，$N^2=\zero$。级数从 $k=2$ 起全为零：
\[
e^{Nx}=E+Nx=\pmat{1&x\\0&1}.
\]
这个例子还说明：即使 $A$ 不能对角化，$e^{Ax}$ 也完全可以有简洁形式——这正是 §4 要系统处理的情形。

\subhdr{1.7 为什么还需要别的算法}

按定义算 $e^{Ax}$ 要算无穷多项，除对角、纯量、幂零等特殊情形外不实用（课件 p6、p17 都强调了这一点）。
后文两条路线分别利用 $A$ 的两种结构：
\begin{itemize}
  \item 有 $n$ 个线性无关特征向量时，$A$ 被对角化，$e^{Ax}$ 化为 $n$ 个标量指数；
  \item 特征向量不够时，Cayley--Hamilton 定理保证高次幂可用低次幂表示，
        于是 $e^{Ax}$ 只需算到 $A^{n-1}$，系数由特征值处的插值条件确定。
\end{itemize}

%=====================================================================
\hdr{§2}{矩阵指数的性质（p7）}
%=====================================================================

\subhdr{2.1 Property 1.1：五条性质及各自的依据}

设 $A,B$ 是 $n\times n$ 矩阵，$r,s,t$ 是数。

\textbf{(1) $e^{A0}=e^{\zero}=E$.} 级数中只有 $k=0$ 项：$A^0x^0/0!=E$，其余项含 $x=0$ 因子。

\textbf{(2) $e^{At}e^{As}=e^{A(t+s)}$.} 两个级数都绝对收敛，可以按 Cauchy 乘积重排：
\[
e^{At}e^{As}
=\sum_{n=0}^{\infty}\sum_{k=0}^{n}
  \frac{A^kt^k}{k!}\,\frac{A^{n-k}s^{n-k}}{(n-k)!}
=\sum_{n=0}^{\infty}\frac{A^n}{n!}\sum_{k=0}^{n}\binom{n}{k}t^ks^{n-k}
=\sum_{n=0}^{\infty}\frac{A^n(t+s)^n}{n!}.
\]
关键一步是 $A^kA^{n-k}=A^n$，这只用到“同一个 $A$ 的幂彼此可交换”，不需要 $A$ 与别的矩阵可交换。

\textbf{(3) 若 $AB=BA$，则 $e^{A+B}=e^Ae^B$.} 此时 $(A+B)^n$ 可按二项式展开：
$(A+B)^n=\sum_{k=0}^n\binom{n}{k}A^kB^{n-k}$，再按 (2) 的乘法重排即得。
条件 $AB=BA$ 不能省；反例见 2.3。

\textbf{(4) $e^{-At}=(e^{At})^{-1}$.} 在 (2) 中取 $s=-t$：
$e^{At}e^{-At}=e^{A(t-t)}=e^{A0}=E$；同理 $e^{-At}e^{At}=E$。这同时给出 $e^{At}$ 可逆。

\textbf{(5) $e^{rtE}=e^{rt}E$.} $rtE$ 是对角元全为 $rt$ 的对角矩阵，由 1.6 例 B 即得。

\subhdr{2.2 Property 1.2：特征向量经过 $e^{A}$ 仍是特征向量}

\callout{
\textbf{Property 1.2（p7）.} 若 $V\in\R^n$ 满足 $AV=\lambda V$（$V\ne0$），则
$e^{A}V=e^{\lambda}V$；更一般地，$e^{Ax}V=e^{\lambda x}V$。}

\textbf{证明.} 由 $AV=\lambda V$ 归纳得 $A^kV=\lambda^kV$（$k=0$ 时 $V=\lambda^0V$）。于是
\[
e^{Ax}V=\sum_{k=0}^{\infty}\frac{A^kx^k}{k!}V
=\sum_{k=0}^{\infty}\frac{\lambda^kx^k}{k!}V=e^{\lambda x}V.
\]
这说明：$A$ 的特征向量方向在 $e^{Ax}$ 作用下方向不变，只被乘上标量因子 $e^{\lambda x}$。
§3 的对角化公式本质上就是把这句话逐列写出来。

\subhdr{2.3 易错点：乘法顺序与 $AB=BA$}

$e^{A+B}=e^{A}e^{B}$ 依赖 $AB=BA$。若两者不可交换，等号一般不成立。取
\[
A=\pmat{0&1\\0&0},\qquad B=\pmat{0&0\\1&0}.
\]
则 $A^2=B^2=\zero$，所以
\[
e^{A}=E+A=\pmat{1&1\\0&1},\qquad
e^{B}=E+B=\pmat{1&0\\1&1},\qquad
e^{A}e^{B}=\pmat{2&1\\1&1}.
\]
而 $A+B=\pmat{0&1\\1&0}$ 满足 $(A+B)^2=E$，故
\[
e^{A+B}=\cosh(1)E+\sinh(1)(A+B)
=\pmat{\cosh1&\sinh1\\\sinh1&\cosh1}
\ne\pmat{2&1\\1&1}.
\]
判断方法只有一条：先算 $AB$ 与 $BA$。相等才能用 (3)；不相等只能用级数定义或其它结构。

另一个常见误区是把 $e^{At}e^{As}$ 与 $e^{A t}e^{B s}$ 混为一谈：前者对同一个 $A$，总能合并；
后者涉及两个矩阵，只有在 $AB=BA$ 时才能把指数相加。

\flowerrule

%=====================================================================
\clearpage
\hdr{§3}{对角化方法：$g=m$ 情形（p8--p15）}
%=====================================================================

\subhdr{3.1 特征方程、代数重数与几何重数}

设 $A$ 是 $n\times n$ 矩阵。取首一特征多项式
\[
\Delta(\lambda)=\det(\lambda E-A)
\]
（与 $\det(A-\lambda E)$ 相差 $(-1)^n$，根相同），它的根是特征值。若 $\lambda_1$ 是 $m_1$ 重根，称 $m_1$ 为它的\textbf{代数重数}；而
\[
g_1=n-\operatorname{rank}(A-\lambda_1E)
\]
是特征子空间 $\ker(A-\lambda_1E)$ 的维数，称为\textbf{几何重数}。总有 $1\le g_1\le m_1$。

课件 p9 把特征值按代数重数排成
$\lambda_1,\dots,\lambda_1,\lambda_{m_1+1},\dots,\lambda_n$（前 $m_1$ 个都是 $\lambda_1$），
并写 $g_1=m_1$。这是“Case 1”的条件。要强调：这个条件必须对\emph{每个}特征值都成立。
判断步骤是：
\begin{enumerate}
  \item 解 $\det(A-\lambda E)=0$，记下每个特征值的代数重数 $m$；
  \item 对每个特征值，对 $A-\lambda E$ 做行变换，数非零行数得秩，算出 $g=n-\operatorname{rank}$；
  \item 若所有特征值都有 $g=m$，则 $A$ 可对角化，本节公式可用；
        只要有一个 $g<m$，特征向量总数不足 $n$，$V$ 不可逆，必须转 §4。
\end{enumerate}

\subhdr{3.2 拼出 $V$ 与 $\Lambda$}

取每个特征值的 $m_j$ 个线性无关特征向量（$g_j=m_j$ 保证取得到），按列拼成
\[
V=(V_1,V_2,\dots,V_n).
\]
由于 $n$ 个列线性无关，$\det V\ne0$。于是
\[
AV=A(V_1,\dots,V_n)=(AV_1,\dots,AV_n)=(\lambda_1V_1,\dots,\lambda_nV_n)=V\Lambda,
\]
其中 $\Lambda=\operatorname{diag}(\lambda_1,\dots,\lambda_n)$。右乘 $V^{-1}$ 得
\begin{equation}\label{eq:diag}
A=V\Lambda V^{-1}.
\end{equation}
用之前请实际验一次 $AV=V\Lambda$（逐列核对 $AV_j=\lambda_jV_j$），这比反复检查 $V^{-1}$ 更省事。

\subhdr{3.3 推导 $e^{Ax}=Ve^{\Lambda x}V^{-1}$}

把 \eqref{eq:diag} 代入定义。因为 $(V\Lambda V^{-1})^k=V\Lambda^kV^{-1}$（中间相邻的
$V^{-1}V$ 两两抵消），
\[
\begin{aligned}
e^{Ax}
&=E+Ax+\frac{x^2}{2!}(V\Lambda V^{-1})(V\Lambda V^{-1})+\cdots
+\frac{x^k}{k!}(V\Lambda V^{-1})^k+\cdots\\
&=E+V\Lambda V^{-1}x+\frac{x^2}{2!}V\Lambda^2V^{-1}+\cdots
+\frac{x^k}{k!}V\Lambda^kV^{-1}+\cdots\\
&=V\left(\sum_{k=0}^{\infty}\frac{\Lambda^kx^k}{k!}\right)V^{-1}
=Ve^{\Lambda x}V^{-1}.
\end{aligned}
\]
第一步把每个 $(\Ax)^k$ 展开；第二步用 $V^{-1}V=E$；第三步把公因子 $V$、$V^{-1}$ 提到级数外
（有限矩阵、绝对收敛，合法），其中 $k=0$ 项是 $E=V\cdot E\cdot V^{-1}$。
\needspace{8\baselineskip}
最后因为 $\Lambda$ 是对角矩阵，
\[
e^{\Lambda x}=\operatorname{diag}\left(e^{\lambda_1x},\dots,e^{\lambda_nx}\right),
\]
于是
\begin{equation}\label{eq:diagExp}
e^{Ax}=V\operatorname{diag}\left(e^{\lambda_1x},\dots,e^{\lambda_nx}\right)V^{-1}.
\end{equation}
这就是课件 p10--p11 的结论。它的使用前提只有一个：$V$ 的列确实是 $n$ 个线性无关的特征向量。

\subhdr{3.4 通解与初值问题}

由 \eqref{eq:diagExp}，若记 $V=(V_1,\dots,V_n)$，则
\[
Y(x)=Ve^{\Lambda x}C
=c_1V_1e^{\lambda_1x}+c_2V_2e^{\lambda_2x}+\cdots+c_nV_ne^{\lambda_nx},
\]
其中 $C=(c_1,\dots,c_n)^T$ 是任意常向量。这就是 $Y'=AY$ 的通解：它是 $n$ 个线性无关解
$V_je^{\lambda_jx}$ 的线性组合。

对初值问题 \eqref{eq:system}，把 $x_0$ 代入：
\[
Y(x)=e^{A(x-x_0)}Y_0=Ve^{\Lambda(x-x_0)}V^{-1}Y_0 .
\]
课件 p12 把 $e^{-\Lambda x_0}V^{-1}Y_0$ 记成 $(c_1,\dots,c_n)^T$，于是
$Y=Ve^{\Lambda x}(c_1,\dots,c_n)^T$。注意这些 $c_j$ 依赖 $x_0$ 与 $Y_0$，不是任意常数；
若要求通解，就让 $C$ 任意。

\subhdr{3.5 课件例题 1（教学示范，p13--p15）}

课件在 p13 明确把它标为“Lecture 14--15 Example 1”的重现，并在 p14--p15 给出完整解答。
本讲义把它作为\textbf{教学示范}完整讲一遍：它同时演示了特征值、特征向量、$V^{-1}$ 的行变换与
对角化公式。此例不作为独立练习。

\textbf{题目.} 求 $A$ 的 $e^{Ax}$ 与 $Y'=AY$ 的通解，其中
\[
A=\pmat{1&-1&1\\1&1&-1\\2&-1&0}.
\]

\textbf{第一步：特征值.} 计算
\[
\det(A-\lambda E)=-(\lambda-2)(\lambda-1)(\lambda+1).
\]
三个特征值 $\lambda_1=2$、$\lambda_2=1$、$\lambda_3=-1$，代数重数都是 1，因此几何重数也必为 1，
满足 $g=m$。

\textbf{第二步：特征向量.} 分别解 $(A-\lambda_jE)V=0$：
\[
\lambda_1=2:\ V_1=\pmat{1\\0\\1},\qquad
\lambda_2=1:\ V_2=\pmat{1\\1\\1},\qquad
\lambda_3=-1:\ V_3=\pmat{-1\\3\\5}.
\]
逐列验算 $AV_j=\lambda_jV_j$：例如
$AV_1=(1+0+1,\;1+0-1,\;2+0+0)^T=(2,0,2)^T=2V_1$，另两个同理。于是
\[
V=\pmat{1&1&-1\\0&1&3\\1&1&5},\qquad
\Lambda=\pmat{2&0&0\\0&1&0\\0&0&-1}.
\]
$\det V=6\ne0$，$V$ 可逆。

\textbf{第三步：求 $V^{-1}$.} 对 $[V\mid E]$ 做行变换：
\[
\left(\begin{array}{ccc|ccc}
1&1&-1&1&0&0\\
0&1&3&0&1&0\\
1&1&5&0&0&1
\end{array}\right)
\xrightarrow{R_3-R_1}
\left(\begin{array}{ccc|ccc}
1&1&-1&1&0&0\\
0&1&3&0&1&0\\
0&0&6&-1&0&1
\end{array}\right)
\]
\clearpage
\noindent\textbf{例 1（续）}\quad 继续对 $[V\mid E]$ 做行变换。
\[
\xrightarrow{R_2-\frac12R_3,\ R_3/6}
\left(\begin{array}{ccc|ccc}
1&1&-1&1&0&0\\
0&1&0&\frac12&1&-\frac12\\
0&0&1&-\frac16&0&\frac16
\end{array}\right)
\xrightarrow{R_1-R_2+R_3}
\left(\begin{array}{ccc|ccc}
1&0&0&\frac13&-1&\frac23\\
0&1&0&\frac12&1&-\frac12\\
0&0&1&-\frac16&0&\frac16
\end{array}\right).
\]
所以
\[
V^{-1}=\pmat{\frac13&-1&\frac23\\[2pt]\frac12&1&-\frac12\\[2pt]-\frac16&0&\frac16}.
\]
\clearpage
\noindent\textbf{例 1（续）}\quad 核对 $V^{-1}$ 后写出 $e^{Ax}$ 与通解；当前
$A=\pmat{1&-1&1\\1&1&-1\\2&-1&0}$。

（课件 p14 的最后一行把 $(1,2)$ 元素写成了 $0$；用 $VV^{-1}=E$ 或代数余子式可验应为 $-1$。
这一处按勘误处理，附录 B 给出核对。若沿用印刷值，$x=0$ 时 $Ve^{\Lambda\cdot0}V^{-1}=VV^{-1}\ne E$。）

\textbf{第四步：写出 $e^{Ax}$ 与通解.} 由 \eqref{eq:diagExp}，
\[
e^{Ax}=V\pmat{e^{2x}&0&0\\0&e^{x}&0\\0&0&e^{-x}}V^{-1}.
\]
乘开后是
\[
e^{Ax}=
\pmat{
\frac13e^{2x}+\frac12e^{x}+\frac16e^{-x} &
-e^{2x}+e^{x} &
\frac23e^{2x}-\frac12e^{x}-\frac16e^{-x}\\[4pt]
\frac12e^{x}-\frac12e^{-x} & e^{x} & -\frac12e^{x}+\frac12e^{-x}\\[4pt]
\frac13e^{2x}+\frac12e^{x}-\frac56e^{-x} &
-e^{2x}+e^{x} &
\frac23e^{2x}-\frac12e^{x}+\frac56e^{-x}}
\]
（这一行是把 $V$、$e^{\Lambda x}$、$V^{-1}$ 三个矩阵实际相乘的结果，可用 §0.4 的三条自检核对）。
于是通解为
\[
Y(x)=Ve^{\Lambda x}C
=c_1\pmat{1\\0\\1}e^{2x}+c_2\pmat{1\\1\\1}e^{x}+c_3\pmat{-1\\3\\5}e^{-x}.
\]
若给定 $Y(x_0)=Y_0$，则 $Y=Ve^{\Lambda(x-x_0)}V^{-1}Y_0$。

\textbf{自检.} (1) $x=0$ 时 $e^{A0}=E$（三个对角元 $1/3+1/2+1/6=1$ 等）；(2) 求导后令 $x=0$ 得
\[
\left.\frac{\dd}{\dd x}e^{Ax}\right|_{x=0}
=\pmat{1&-1&1\\1&1&-1\\2&-1&0}=A;
\]
(3) 每一列都满足 $Y'=AY$（因为 $e^{Ax}$ 满足矩阵方程，见 1.4）。

\needspace{14\baselineskip}
\subhdr{3.6 条件不满足时会发生什么}

若某个特征值 $g<m$，线性无关特征向量不足 $n$ 个，$V$ 不可逆，\eqref{eq:diagExp} 就不能用。
最简例子是
\[
A=\pmat{\lambda&1\\0&\lambda}.
\]
它只有一个特征值 $\lambda$（代数重数 2），而
$g=2-\operatorname{rank}(A-\lambda E)=2-1=1<2$。这时仍可算出
\[
e^{Ax}=e^{\lambda x}\pmat{1&x\\0&1},
\]
但无法用“$n$ 个特征向量对角化”的方式得到。§4 的 Cayley--Hamilton 方法不要求可对角化，
正是为这种情形准备的。

\flowerrule

%=====================================================================
\hdr{§4}{Cayley--Hamilton 方法（p16--p38）}
%=====================================================================

\subhdr{4.1 动机与 Cayley--Hamilton 定理}

对角化要求“特征向量够用”。本节的方法只依赖特征多项式，对任何方阵都成立，因此能处理重特征值。
先回顾定理（课件 p17，Theorem 3.1）。

\callout{
\textbf{Cayley--Hamilton 定理.} 设 $A$ 是 $n\times n$ 矩阵。取首一特征多项式
\[
p(\lambda)=\det(\lambda E-A)=\lambda^n+c_{n-1}\lambda^{n-1}+\cdots+c_1\lambda+c_0 .
\]
则
\[
p(A)=A^n+c_{n-1}A^{n-1}+\cdots+c_1A+c_0E=\zero .
\]}

课件把特征多项式写成 $\det(A-\lambda E)=b_n\lambda^n+\cdots+b_0$。它与首一版本
$p(\lambda)=\det(\lambda E-A)$ 只差 $(-1)^n$（即 $b_n=(-1)^n$），根相同，“在 $A$ 处取值为零矩阵”
的结论也相同。本讲义在需要“首一”或做除法时一律用 $p$（后文记作 $\Delta$）；只求特征值时写
$\det(A-\lambda E)=0$ 也可以。证明超出本讲范围，但定理的用法必须会：
它把 $A^n$ 表示成 $E,A,\dots,A^{n-1}$ 的线性组合，从而任何高次幂都能降阶。

\textbf{教学小例.} 取 $A=\pmat{3&1\\1&2}$，特征多项式
$\det(\lambda E-A)=\lambda^2-5\lambda+5$。定理给出 $A^2-5A+5E=\zero$，即
\[
A^2=5A-5E.
\]
于是 $A^3=5A^2-5A=5(5A-5E)-5A=20A-25E$，任意高次幂都能压到一次。这个矩阵马上会在例题 2 中再用到。

\needspace{12\baselineskip}
\subhdr{4.2 从无穷级数到 $n$ 项多项式（Theorem 3.2，p18）}

\callout{
\textbf{Theorem 3.2（p18）.} 设 $A$ 是 $n\times n$ 矩阵，则存在只依赖 $x$ 的函数
$\alpha_0(x),\dots,\alpha_{n-1}(x)$，使下式成立：}
\begin{equation}\label{eq:reduce}
e^{Ax}=\alpha_{n-1}(x)A^{n-1}x^{n-1}+\cdots+\alpha_1(x)Ax+\alpha_0(x)E
=\sum_{k=0}^{n-1}\alpha_k(x)(Ax)^k .
\end{equation}

\textbf{为什么只有 $n$ 项.} 固定一个 $x$，把 $\Ax$ 看成一个普通矩阵。它的特征多项式
$\Delta_{\Ax}(s)=\det(sE-\Ax)$（首一；若用 $\det(\Ax-sE)$ 则差 $(-1)^n$）满足矩阵恒等式
$\Delta_{\Ax}(\Ax)=\zero$，其系数是 $x$ 的多项式。由 Cayley--Hamilton，
$(\Ax)^n$ 可写成 $E,\Ax,\dots,(\Ax)^{n-1}$ 的线性组合；反复代入，任何 $(\Ax)^k$（$k\ge n$）
都可降到次数 $\le n-1$。对 $e^{Ax}$ 的部分和逐项降阶，再令项数趋于无穷，降阶后的系数收敛，
得到 \eqref{eq:reduce}。严格说，这一步用到矩阵级数在任意有限区间上的一致收敛；
本讲只需掌握“为什么只有 $n$ 项”这一机制。这里的要点是：系数 $\alpha_k$ 一般\emph{依赖 $x$}，因为 $\Ax$ 本身含 $x$；
不能把它们当成常数。

\subhdr{4.3 多项式降阶与长除法（p19）}

若 $p$ 是任一多项式，用首一特征多项式 $\Delta(s)=\det(sE-A)$ 做带余除法（普通有限多项式长除法）：
\[
p(s)=q(s)\Delta(s)+r(s),\qquad \deg r\le n-1 .
\]
在 $s=A$ 处取值，由 $\Delta(A)=\zero$ 得
\begin{equation}\label{eq:remainder}
p(A)=r(A).
\end{equation}
商 $q$ 由长除法得到，余式 $r$ 的次数不超过 $n-1$。注意：这里不是“近似”，而是矩阵恒等式；
不要漏掉 $E$（常数项对应 $cE$，不是数 $c$）。

\clearpage
\subhdr{4.4 课件例题 2（教学示范，p20）}

\textbf{题目.} 设 $A=\pmat{3&1\\1&2}$，把 $p(A)=A^4+3A^3+2A^2+A+E$ 降阶。

\textbf{解.} 首一特征多项式 $\Delta(s)=\det(sE-A)=s^2-5s+5$（$n=2$，与 $\det(A-sE)$ 相同）。做长除法：
\[
\frac{s^4+3s^3+2s^2+s+1}{s^2-5s+5}
=s^2+8s+37+\frac{146s-184}{s^2-5s+5}.
\]
核对乘法：
$(s^2+8s+37)(s^2-5s+5)=s^4+3s^3+2s^2-145s+185$，加上余式 $146s-184$ 正是 $p(s)$。
于是
\[
p(s)=(s^2+8s+37)\Delta(s)+(146s-184),\qquad r(s)=146s-184,
\]
从而
\[
p(A)=146A-184E
=146\pmat{3&1\\1&2}-184\pmat{1&0\\0&1}
=\pmat{254&146\\146&108}.
\]
\textbf{自检.} 直接按定义算 $p(A)$ 也能得到同一矩阵（计算量更大）；至少检查迹与行列式：
$\operatorname{tr}p(A)=362$，与 $146\operatorname{tr}A-184\cdot2=146\cdot5-368=362$ 一致。

\subhdr{4.5 解析函数、特征多项式与插值方程组（p21--p22）}

先固定 $A$，记它的首一特征多项式为
\[
\Delta(s)=\det(sE-A)=s^n+c_{n-1}s^{n-1}+\cdots+c_0 .
\]
（若用 $\det(A-sE)$，则整体差 $(-1)^n$，根相同；需要“首一”或做除法时一律用 $\Delta$。）

\textbf{多项式情形：带余除法.} 对任一多项式 $P(s)$，普通有限多项式长除法给出
\[
P(s)=\Delta(s)q(s)+r(s),\qquad \deg r\le n-1 .
\]
设 $\lambda_i$ 是 $\Delta$ 的 $m_i$ 重根（$\sum_i m_i=n$）。因为 $\Delta$ 在 $\lambda_i$ 处有
$m_i$ 阶零点，$P-r=\Delta q$ 也有 $m_i$ 阶零点，所以
\begin{equation}\label{eq:interp}
P^{(j)}(\lambda_i)=r^{(j)}(\lambda_i),\qquad j=0,1,\dots,m_i-1 .
\end{equation}

\textbf{矩阵指数情形：先定 $r$，再验证整除.} 对 $f(s)=e^s$，它不是多项式，不能把上面的长除法
无条件照搬成“任何幂级数都能做长除法”。正确的做法是反过来：先取次数 $\le n-1$ 的多项式 $r$，
使 \eqref{eq:interp} 对 $f$ 成立。这是 $n$ 个条件、$n$ 个系数 $\alpha_0,\dots,\alpha_{n-1}$ 的
Hermite 插值问题，存在且唯一。再令
\[
q(s)=\frac{f(s)-r(s)}{\Delta(s)} .
\]
由于 $f-r$ 在每个 $\lambda_i$ 处有 $m_i$ 阶零点，而 $\Delta$ 的零点阶数正好是 $m_i$，所以 $q$ 在
$\lambda_i$ 处解析；在其它点也解析。因此 $q$ 是整函数，且
\begin{equation}\label{eq:div}
f(s)=\Delta(s)q(s)+r(s)
\end{equation}
对全部 $s$ 成立。这一步不是“把幂级数直接长除”，而是用插值条件保证分子被分母整除；
对有限多项式，它就是普通带余除法。

等价地说，$r$ 是 $e^s$ 在 $A$ 的谱上的 Hermite 插值多项式。若 $n$ 个特征值互不相同
（都取 $m_i=1$），\eqref{eq:interp} 就是 $n$ 个方程、$n$ 个未知量的 Vandermonde 线性方程组；
解出后 $f(A)=r(A)=\sum_{k=0}^{n-1}\alpha_kA^k$。至于为什么这些条件能推出 $e^{Ax}=r(\Ax)$，
见 §4.9 的 Jordan 块论证。

对矩阵指数，把 $A$ 换成 $\Ax$，特征多项式换成 $\Delta_x(s)=\det(sE-\Ax)$，特征值换成
$\lambda_i(x)=x\mu_i$，其中 $\mu_i$ 是 $A$ 的特征值；当 $x\ne0$ 时重数 $m_i$ 不变。于是
\begin{equation}\label{eq:eigenEq}
e^{\lambda_i}=r^{(j)}(\lambda_i),\qquad j=0,1,\dots,m_i-1,\qquad
r(\lambda)=\sum_{k=0}^{n-1}\alpha_k\lambda^k .
\end{equation}
请再次注意：$\alpha_k$ 是 $x$ 的函数，而 $\lambda$ 是插值用的\emph{标量变量}；
对 $\lambda$ 求导时把 $\alpha_k$ 当常数。

\subhdr{4.5.1 适用范围：$A$ 的固定特征值与 $\Ax$ 的伸缩特征值}

这里必须区分两层。$A$ 的特征值 $\mu_1,\dots,\mu_s$ 是固定的，代数重数为 $m_1,\dots,m_s$
（$m_1+\cdots+m_s=n$）；$\Ax$ 的特征值是 $\lambda_i(x)=x\mu_i$。
\begin{itemize}
  \item 当 $x\ne0$ 时，$\lambda_i(x)$ 与 $\mu_i$ 一一对应，重数相同；这时才可以把 $x\mu_i$ 当作
        $n$ 个不同的节点，或在推导中除以 $x$、$x^2$。
  \item 当 $x=0$ 时，$\Ax=0$，所有 $\lambda_i(0)=0$ 合并成 $0$，代数重数为 $n$，插值方程退化
        （都变成同一点 $0$ 的导数条件）。这时结论由极限得到：例子中 $\alpha_k(x)$ 的表达式
        可能在 $0$ 处分母为零，但极限存在；把 $\alpha_k(0)$ 定义为该极限，$e^{A\cdot0}=E$ 仍成立。
        这不是“补一个特例”，而是因为 $e^{Ax}$ 与 $\alpha_k(x)$ 都是 $x$ 的连续（解析）函数，
        等式在 $x\ne0$ 成立后可在 $x=0$ 取极限。
\end{itemize}

\subhdr{4.6 课件例题 3（教学示范，p23）}

\textbf{题目.} 说明当 $A$ 是 $2\times2$ 或 $3\times3$ 时，$e^{Ax}$ 的待定形式。

\textbf{解.} 当 $n=2$ 时，$\deg r\le1$，所以
\[
e^{Ax}=\alpha_1(x)\Ax+\alpha_0(x)E .
\]
当 $n=3$ 时，$\deg r\le2$，所以
\[
e^{Ax}=\alpha_2(x)A^2x^2+\alpha_1(x)\Ax+\alpha_0(x)E .
\]
一般地就是 \eqref{eq:reduce}。这里的关键不是“猜测”，而是 4.2 的降阶结论：$\Ax$ 的 $n$ 次及以上幂
都能用前 $n$ 个幂表示，因此只需 $n$ 个待定函数。课件的 Remark 强调 $\alpha_i(x)$ 由
$\Ax$ 的特征值所给出的方程组确定；下一小节起就是具体做法。

\subhdr{4.7 课件例题 4（教学示范，p24--p25）：两个不同实特征值}

\textbf{题目.} 求 $A=\pmat{1&1\\9&1}$ 的 $e^{Ax}$，并解 $Y'=AY$。

\textbf{第一步：写出待定形式.} $n=2$，
\[
e^{Ax}=\alpha_1\Ax+\alpha_0E
=\pmat{\alpha_1x+\alpha_0&\alpha_1x\\9\alpha_1x&\alpha_1x+\alpha_0}.
\]
注意矩阵里出现的是 $\alpha_1x$，所以后面直接解 $\alpha_1x$ 与 $\alpha_0$ 更省事。

\clearpage
\noindent\textbf{例 4（续）}\quad 当前 $A=\pmat{1&1\\9&1}$，$e^{Ax}=\alpha_1\Ax+\alpha_0E$；接着求
$\Ax$ 的特征值。

\textbf{第二步：特征值.}
\[
\det(\Ax-\lambda E)
=\begin{vmatrix}x-\lambda&x\\9x&x-\lambda\end{vmatrix}
=(x-\lambda)^2-9x^2=(\lambda-4x)(\lambda+2x),
\]
故对 $x\ne0$，$\Ax$ 的特征值是 $\lambda_1=4x$、$\lambda_2=-2x$，互不相同；对应的 $A$ 的固定特征值是
$\mu_1=4$、$\mu_2=-2$，它们本身也互不相同。$x=0$ 时两个特征值都变成 $0$，插值方程退化，
按 §4.5.1 取极限（最终闭式在 $x=0$ 给出 $E$）。

\textbf{第三步：列方程.} 由 $e^{\lambda_i}=r(\lambda_i)$，$r(\lambda)=\alpha_1\lambda+\alpha_0$：
\[
e^{4x}=4x\alpha_1+\alpha_0,\qquad
e^{-2x}=-2x\alpha_1+\alpha_0 .
\]
写成矩阵：
\[
\pmat{4&1\\-2&1}\pmat{x\alpha_1\\\alpha_0}=\pmat{e^{4x}\\e^{-2x}}.
\]

\textbf{第四步：解方程.}
两式相减得 $6x\alpha_1=e^{4x}-e^{-2x}$，即
\[
x\alpha_1=\frac16\left(e^{4x}-e^{-2x}\right),\qquad
\alpha_0=\frac13\left(e^{4x}+2e^{-2x}\right).
\]
代回第一步：
\[
e^{Ax}=\frac16\pmat{
3e^{4x}+3e^{-2x} & e^{4x}-e^{-2x}\\
9e^{4x}-9e^{-2x} & 3e^{4x}+3e^{-2x}} .
\]
\textbf{第五步：通解.}
\[
Y=e^{Ax}C
=\tilde c_1e^{4x}\pmat{1\\3}+\tilde c_2e^{-2x}\pmat{1\\-3},
\qquad
\tilde c_1=\frac{3c_1+c_2}{6},\quad
\tilde c_2=\frac{3c_1-c_2}{6}.
\]
这里 $\tilde c_1,\tilde c_2$ 中的因子 $1/6$ 不能漏。课件 p25 把 $\tilde c_1$ 写成 $3c_1+c_2$、
$\tilde c_2$ 写成 $3c_1-c_2$，相当于漏掉公共因子 $1/6$；用 $x=0$ 时 $Y(0)=(c_1,c_2)^T$
即可验出。附录 B 列有这一处勘误。注意这里固定的是同一个矩阵等式两边的常数；若只关心通解族，
任意常数可以重新命名，$1/6$ 可被吸收到新的常数里，那是另一回事，不能用来解释这处等式错误。

\textbf{自检.} (1) $x=0$ 时 $e^{A0}=\frac16\pmat{6&0\\0&6}=E$；
(2)逐项求导后令 $x=0$：
\[
\left.\frac{\dd}{\dd x}e^{Ax}\right|_{x=0}
=\frac16\pmat{12-6&4+2\\36+18&-6+12}
=\frac16\pmat{6&6\\54&6}
=\pmat{1&1\\9&1}=A;
\]
(3) 每一列都满足 $Y'=AY$（矩阵解的性质）。

\subhdr{4.8 课件例题 5（教学示范，p26--p27）：共轭复特征值}

\textbf{题目.} 解 $Y'=AY$，其中 $A=\pmat{0&1\\-1&0}$。

\clearpage
\noindent\textbf{例 5（续）}\quad 当前 $A=\pmat{0&1\\-1&0}$，求 $e^{Ax}$ 与 $Y'=AY$ 的通解。

\textbf{解.} $n=2$，
\[
e^{Ax}=\alpha_1\Ax+\alpha_0E
=\pmat{\alpha_0&\alpha_1x\\-\alpha_1x&\alpha_0},\qquad
r(\lambda)=\alpha_1\lambda+\alpha_0 .
\]
特征多项式
\[
\det(\Ax-\lambda E)=\begin{vmatrix}-\lambda&x\\-x&-\lambda\end{vmatrix}=\lambda^2+x^2,
\]
对 $x\ne0$，特征值 $\lambda_1=\ii x$、$\lambda_2=-\ii x$ 互不相同；对应的 $A$ 的固定特征值是
$\pm\ii$，也互不相同。$x=0$ 时两者都变成 $0$，按 §4.5.1 取极限（最终闭式在 $x=0$ 给出 $E$）。
由 $e^{\lambda_i}=r(\lambda_i)$：
\[
e^{\ii x}=\ii x\alpha_1+\alpha_0,\qquad
e^{-\ii x}=-\ii x\alpha_1+\alpha_0 .
\]
两式相减、相加并用 Euler 公式：
\[
\alpha_1x=\frac{e^{\ii x}-e^{-\ii x}}{2\ii}=\sin x,\qquad
\alpha_0=\frac{e^{\ii x}+e^{-\ii x}}{2}=\cos x .
\]
于是
\[
e^{Ax}=\pmat{\cos x&\sin x\\-\sin x&\cos x}.
\]
通解为
\[
Y(x)=e^{Ax}C
=c_1\pmat{\cos x\\-\sin x}+c_2\pmat{\sin x\\\cos x}.
\]
课件 p27 把第二个向量写成 $(-\sin x,\cos x)^T$，与它自己上一行的矩阵
$\pmat{\cos x&\sin x\\-\sin x&\cos x}$ 的第二列 $( \sin x,\cos x)^T$ 矛盾。
用求导自检：正确的 $Y'(0)=c_1(0,-1)^T+c_2(1,0)^T=(c_2,-c_1)^T$，而
$AY(0)=A(c_1,c_2)^T=(c_2,-c_1)^T$，一致；若取课件的 $(-\sin x,\cos x)^T$，则
$Y'(0)=(-c_2,-c_1)^T$，不满足方程。故按勘误取正号。

\textbf{自检.} (1) $x=0$ 得 $E$；(2)导数为
$\pmat{-\sin x&\cos x\\-\cos x&-\sin x}=Ae^{Ax}$；(3) $\det e^{Ax}=1\ne0$，
两列线性无关（下一节会用到）。

\subhdr{4.9 重特征值：方程不够，需要导数条件（p28--p29）}

若 $\lambda_i=\lambda_j$（$i\ne j$），则 \eqref{eq:eigenEq} 中两个特征值的零阶方程完全相同：
$e^{\lambda_i}=\sum_k\alpha_k\lambda_i^k$。$n$ 个未知量却只有少于 $n$ 个独立方程，解不出来。
解决办法不是“把只在一点成立的等式 $e^\lambda=r(\lambda)$ 直接对 $\lambda$ 求导”，而是利用
§4.5 的整除关系 \eqref{eq:div}：
\[
f(s)-r(s)=\Delta_x(s)q(s),\qquad f(s)=e^s .
\]
\needspace{10\baselineskip}
设 $\lambda_i$ 是 $\Delta_x$ 的 $m_i$ 重根。因为 $\Delta_x$ 在 $\lambda_i$ 处有 $m_i$ 阶零点，右边
也有 $m_i$ 阶零点，所以左边在 $\lambda_i$ 处有 $m_i$ 阶零点。于是
\begin{equation}\label{eq:derivCond}
e^{\lambda_i}=r(\lambda_i),\quad
e^{\lambda_i}=r'(\lambda_i),\quad
\dots,\quad
e^{\lambda_i}=r^{(m_i-1)}(\lambda_i).
\end{equation}
（这里 $\lambda_i$ 是 $\Ax$ 的特征值；$\alpha_k$ 只依赖 $x$，对 $\lambda$ 求导时视为常数。）
这些条件与其它特征值的条件合起来共 $n$ 个，唯一确定 $r$。

\textbf{为什么这些条件足够：Jordan 块.} 还差最后一步：由 \eqref{eq:derivCond} 定出的 $r$ 确实满足
$e^{Ax}=r(\Ax)$。把 $\Ax$ 化成 Jordan 形。对特征值 $\lambda$ 的一个 $\ell$ 阶 Jordan 块
$J=\lambda E+N$（$N^\ell=0$），
\[
e^{J}=e^{\lambda}\sum_{j=0}^{\ell-1}\frac{N^j}{j!},
\qquad
r(J)=\sum_{j=0}^{\ell-1}\frac{r^{(j)}(\lambda)}{j!}N^j
\]
（后一式是 $r$ 在 $\lambda$ 处的 Taylor 展开；$N^\ell=0$ 使 $j\ge\ell$ 的项消失）。Jordan 块的大小
$\ell$ 不超过该特征值的代数重数 $m$。由 \eqref{eq:derivCond}，$e^{\lambda}=r^{(j)}(\lambda)$ 对
$j=0,\dots,m-1$ 成立，当然也对 $j=0,\dots,\ell-1$ 成立。逐项比较 $N^j$ 的系数得 $e^{J}=r(J)$。
$\Ax$ 相似于这些块的直和；若 $\Ax=P\bigl(\oplus J\bigr)P^{-1}$，则
$e^{Ax}=P\bigl(\oplus e^{J}\bigr)P^{-1}=P\bigl(\oplus r(J)\bigr)P^{-1}=r(\Ax)$。故
$e^{Ax}=r(\Ax)$。

\textbf{块大小与代数重数要分清.} 上面的论证只用到“块大小 $\ell\le m$”。如果某个特征值的最大块
只有 $\ell=1$（例如 $A=\mu E$ 是纯量矩阵），那么实际上只需要 $j=0$ 一个条件；我们仍按 $m$ 加了
$j=1,\dots,m-1$ 的导数条件，这些条件不会改变 $r(\Ax)$（因为 $\Ax=\mu xE$，
$r(\Ax)=r(\mu x)E=e^{\mu x}E$），只是把 $r$ 本身选得不同。所以“总代数重数为 $m$”并不推出
“存在 $m$ 阶 Jordan 块”；我们也不需要这个更强的结论，只需要块大小不超过 $m$。

\textbf{$2\times2$ 的具体核对.} 取 $A=\pmat{\mu&1\\0&\mu}$，$N=A-\mu E=\pmat{0&1\\0&0}$，$N^2=\zero$。
于是
\[
e^{Ax}=e^{\mu x}e^{Nx}=e^{\mu x}(E+Nx)=e^{\mu x}\pmat{1&x\\0&1}.
\]
另一方面，$n=2$，$e^{Ax}=\alpha_1\Ax+\alpha_0E$，$\Ax$ 的特征值是 $\mu x$（二重），
$\Delta_x(s)=(s-\mu x)^2$。由 \eqref{eq:derivCond}：
\[
e^{\mu x}=r(\mu x)=\alpha_1\mu x+\alpha_0,\qquad
e^{\mu x}=r'(\mu x)=\alpha_1 .
\]
第二个条件给出 $\alpha_1=e^{\mu x}$，代回得 $\alpha_0=e^{\mu x}(1-\mu x)$，于是
\[
e^{Ax}=e^{\mu x}\Ax+e^{\mu x}(1-\mu x)E
=e^{\mu x}\pmat{1&x\\0&1},
\]
与直接计算一致。这里导数条件是从 $\Delta_x$ 的二重零点得到的，不是对
$e^{\mu x}=r(\mu x)$ 逐点求导。

\callout{
\textbf{Theorem 3.3（插值多项式，完整前提；p29）.} 设 $A$ 是 $n\times n$ 矩阵，$f(s)=e^s$。
记 $\Delta_x(s)=\det(sE-\Ax)$ 为首一特征多项式，其不同根为 $\lambda_1,\dots,\lambda_s$，代数重数为
$m_1,\dots,m_s$（$\sum_i m_i=n$）。令 $r$ 是满足
\[
f^{(j)}(\lambda_i)=r^{(j)}(\lambda_i),\qquad
j=0,1,\dots,m_i-1,\quad i=1,\dots,s
\]
的唯一一个次数 $\le n-1$ 的多项式（Hermite 插值多项式；等价地，$f-r$ 被 $\Delta_x$ 整除）。
则
\[
e^{Ax}=r(\Ax)=\sum_{k=0}^{n-1}\alpha_k(\Ax)^k,
\]
而插值条件就是 $e^{\lambda_i}=r(\lambda_i),e^{\lambda_i}=r'(\lambda_i),\dots,
e^{\lambda_i}=r^{(m_i-1)}(\lambda_i)$。\emph{注意：结论是对这个特定的 $r$ 说的；任意一个次数
$\le n-1$ 的多项式并不自动满足它。}
}

\subhdr{4.10 课件例题 6（教学示范，p30--p31）：三重与单重特征值}

\textbf{题目.} 设 $A$ 是 $4\times4$ 矩阵，$\Ax$ 的特征值为 $\lambda=5x$（三重）与 $\lambda=2x$（一重）。
写出确定 $\alpha_0,\dots,\alpha_3$ 的方程组。

\textbf{解.} $n=4$，
\[
e^{Ax}=\alpha_3A^3x^3+\alpha_2A^2x^2+\alpha_1\Ax+\alpha_0E,\qquad
r(\lambda)=\alpha_3\lambda^3+\alpha_2\lambda^2+\alpha_1\lambda+\alpha_0 .
\]
求导：
\[
r'(\lambda)=3\alpha_3\lambda^2+2\alpha_2\lambda+\alpha_1,\qquad
r''(\lambda)=6\alpha_3\lambda+2\alpha_2 .
\]
对三重特征值 $5x$，用 $e^{5x}=r(5x)=r'(5x)=r''(5x)$：
\[
\begin{aligned}
e^{5x}&=\alpha_3(5x)^3+\alpha_2(5x)^2+\alpha_1(5x)+\alpha_0,\\
e^{5x}&=3\alpha_3(5x)^2+2\alpha_2(5x)+\alpha_1,\\
e^{5x}&=6\alpha_3(5x)+2\alpha_2 .
\end{aligned}
\]
对单重特征值 $2x$，只用 $e^{2x}=r(2x)$：
\[
e^{2x}=\alpha_3(2x)^3+\alpha_2(2x)^2+\alpha_1(2x)+\alpha_0 .
\]
四个方程、四个未知量，系数矩阵的元素是 $x$ 的多项式。

\textbf{补充：把方程组解完（课件在此停住；以下是讲义补写）.} 从最高阶导数方程开始回代：
\[
\alpha_2=\frac{e^{5x}}2-15x\alpha_3,\qquad
\alpha_1=e^{5x}(1-5x)+75x^2\alpha_3,\qquad
\alpha_0=e^{5x}\left(1-5x+\frac{25}{2}x^2\right)-125x^3\alpha_3 .
\]
把这三式代入 $e^{2x}=r(2x)$，整理得
\[
27x^3\alpha_3=e^{5x}\left(1-3x+\frac92x^2\right)-e^{2x},
\]
即
\[
\alpha_3=\frac{e^{5x}(2-6x+9x^2)-2e^{2x}}{54x^3},\qquad
\alpha_2=\frac{e^{5x}(-18x^2+15x-5)+5e^{2x}}{9x^2},
\]
\[
\alpha_1=\frac{e^{5x}(135x^2-132x+50)-50e^{2x}}{18x},\qquad
\alpha_0=\frac{e^{5x}(-225x^2+240x-98)+125e^{2x}}{27}.
\]
\clearpage
\noindent\textbf{例 6（补充，续）}\quad 上页给出 $\alpha_0,\dots,\alpha_3$ 的闭式；下面说明 $x=0$ 的处理。

上面的闭式含 $x$ 的分母，只在 $x\ne0$ 时由上述消元直接得到；$x=0$ 时两个特征值 $5x,2x$ 合并为
$0$（代数重数 $4$），按 §4.5.1 取极限：$\alpha_3\to1/6$、$\alpha_2\to1/2$、$\alpha_1\to1$、
$\alpha_0\to1$，恰好是 $e^{\lambda}$ 在 $\lambda=0$ 处 Taylor 展开的前四项。这个极限检查说明解在
$x=0$ 处没有矛盾。

\subhdr{4.11 求 $e^{Ax}$ 的四步总结（p32）}

\begin{enumerate}
  \item \textbf{求 $\Ax$ 的特征值.} 解 $\det(\lambda E-\Ax)=0$（或 $\det(\Ax-\lambda E)=0$，根相同），把每个特征值按代数重数列出：
        $(\lambda-\lambda_1)^{m_1}\cdots(\lambda-\lambda_s)^{m_s}$，$m_1+\cdots+m_s=n$。
        （注意是 $\Ax$ 的特征值，不是 $A$ 的；它们相差因子 $x$。）
  \item \textbf{写待定形式.} 令
        $e^{Ax}=\alpha_{n-1}(\Ax)^{n-1}+\cdots+\alpha_1\Ax+\alpha_0E$，
        $r(\lambda)=\alpha_{n-1}\lambda^{n-1}+\cdots+\alpha_0$。
  \item \textbf{按重数列方程.} 对每个特征值 $\lambda_i$（重数 $m_i$）列
        \[
        e^{\lambda_i}=r(\lambda_i),\ r'(\lambda_i),\ \dots,\ r^{(m_i-1)}(\lambda_i),
        \]
        一共 $n$ 个方程；解出 $\alpha_0,\dots,\alpha_{n-1}$（它们是 $x$ 的函数）。
  \item \textbf{代回.} 把解出的 $\alpha_k$ 代回第 2 步的待定式，得到 $e^{Ax}$。
\end{enumerate}
最后务必用 §0.4 的三条自检。尤其 $x=0$ 与“求导在 $0$ 处等于 $A$”两条，能抓住常见的符号、
系数与特征向量错误；但它们只是局部检查，通过这两条并不保证整个矩阵函数正确，最后仍要逐列
代回方程（第三条）。

\subhdr{4.12 课件例题 7（教学示范，p33--p35）：三重特征值}

\textbf{题目.} 求
\[
A=\pmat{0&-4&0\\1&-4&0\\1&-2&-2}
\]
的 $e^{Ax}$，并写出 $Y'=AY$ 的解。

\textbf{第一步：$A^2$ 与特征值.}
\[
A^2=\pmat{-4&16&0\\-4&12&0\\-4&8&4}.
\]
特征多项式 $\det(A-\lambda E)=-(\lambda+2)^3$，故 $A$ 只有一个特征值 $\mu=-2$（代数重数 3）；
对 $x\ne0$，$\Ax$ 的特征值是 $\lambda_1=-2x$（三重），$x=0$ 时合并为 $0$（三重），
按 §4.5.1 取极限。因此
\[
e^{Ax}=\alpha_2A^2x^2+\alpha_1\Ax+\alpha_0E,\qquad
r(\lambda)=\alpha_2\lambda^2+\alpha_1\lambda+\alpha_0 .
\]

\textbf{第二步：列方程.} 由 $e^{\lambda}=r(\lambda)$ 及其一、二阶导数在 $\lambda=-2x$ 处取值
（导数条件来自 $\Delta_x$ 的三重零点，见 §4.9）：
\[
e^{-2x}=4x^2\alpha_2-2x\alpha_1+\alpha_0,\qquad
e^{-2x}=-4x\alpha_2+\alpha_1,\qquad
e^{-2x}=2\alpha_2 .
\]
\needspace{10\baselineskip}
写成矩阵（第二、三行乘了 $x$、$x^2$，目的是让未知量统一为 $x^2\alpha_2,x\alpha_1,\alpha_0$；
这与把 $r'(\lambda)$、$r''(\lambda)$ 的方程分别乘 $x$、$x^2$ 等价）：
\[
\pmat{4&-2&1\\-4&1&0\\2&0&0}\pmat{x^2\alpha_2\\x\alpha_1\\\alpha_0}
=\pmat{e^{-2x}\\xe^{-2x}\\x^2e^{-2x}} .
\]

\textbf{第三步：解方程.}
由第三个方程 $\alpha_2=\frac12e^{-2x}$；代入第二个：
\[
x\alpha_1=xe^{-2x}+4x^2\alpha_2=(1+2x)xe^{-2x};
\]
再代入第一个：
\[
\alpha_0=e^{-2x}+2x\alpha_1-4x^2\alpha_2=(1+2x+2x^2)e^{-2x}.
\]
\textbf{第四步：写出 $e^{Ax}$.}
\[
e^{Ax}=\alpha_2x^2A^2+\alpha_1x A+\alpha_0E
=e^{-2x}\pmat{2x+1&-4x&0\\x&-2x+1&0\\x&-2x&1}.
\]
\clearpage
\noindent\textbf{例 7（续）}\quad 当前 $A=\pmat{0&-4&0\\1&-4&0\\1&-2&-2}$，$e^{Ax}$ 见上一页；下面写通解。

\textbf{第五步：通解.} 设 $Y_0=(c_1,c_2,c_3)^T$，则
\[
Y(x)=e^{Ax}Y_0
=e^{-2x}\pmat{(2x+1)c_1-4xc_2\\xc_1+(-2x+1)c_2\\xc_1-2xc_2+c_3}.
\]
把它按“多项式因子 $\times$ 指数”分组：
\[
Y=(c_1-2c_2)e^{-2x}\left[x\pmat{2\\1\\1}+\pmat{1\\0\\0}\right]
+c_2e^{-2x}\pmat{2\\1\\1}
+(c_3-c_2)e^{-2x}\pmat{0\\0\\1}.
\]
课件 p35 把最后一个系数写成 $\tilde c_3=c_3-2c_2$；用 $x=0$ 核对：
$Y(0)$ 的第三个分量应为 $c_3$，而按印刷值会得到 $c_3-c_2$。故应为
$\tilde c_3=c_3-c_2$。附录 B 有这条勘误。

\textbf{自检.} (1) $x=0$ 时 $e^{A0}=E$；(2) 求导在 $0$ 处等于 $A$（可直接对 $e^{-2x}$ 前的矩阵求导验证）；
(3) 三列均满足 $Y'=AY$，因为 $e^{Ax}$ 是矩阵解。

\subhdr{4.13 课件例题 8（教学示范，p36--p38）：$0$ 与 $\pm\ii x$}

\textbf{题目.} 求
\[
A=\pmat{0&1&0\\0&-2&-5\\0&1&2}
\]
的 $e^{Ax}$。

\textbf{第一步：$A^2$ 与特征值.}
\[
A^2=\pmat{0&-2&-5\\0&-1&0\\0&0&-1}.
\]
特征多项式 $\det(A-\lambda E)=-\lambda(\lambda^2+1)$，故 $A$ 的特征值是 $0,\pm\ii$。
对 $x\ne0$，$\Ax$ 的特征值是 $\lambda_1=0$、$\lambda_2=\ii x$、$\lambda_3=-\ii x$，互不相同；
$x=0$ 时三者合并为 $0$（三重），按 §4.5.1 取极限。于是 $n=3$，
\[
e^{Ax}=\alpha_2A^2x^2+\alpha_1\Ax+\alpha_0E,\qquad
r(\lambda)=\alpha_2\lambda^2+\alpha_1\lambda+\alpha_0 .
\]
直接展开待定式可得
\[
e^{Ax}=\pmat{
\alpha_0 & -2\alpha_2x^2+\alpha_1x & -5\alpha_2x^2\\
0 & -\alpha_2x^2-2\alpha_1x+\alpha_0 & -5\alpha_1x\\
0 & \alpha_1x & -\alpha_2x^2+2\alpha_1x+\alpha_0}.
\]
\needspace{12\baselineskip}
\noindent\textbf{例 8（续）}\quad 当前 $A=\pmat{0&1&0\\0&-2&-5\\0&1&2}$；上页矩阵的 $(3,2)$ 元素按勘误取正号，接着列方程。

（课件 p36 把 $(3,2)$ 元素印成 $-\alpha_1x$；按 $\Ax$ 与 $A^2$ 逐项乘加应为 $+\alpha_1x$。
p38 的最终结果也取正号，故按勘误处理。）

\textbf{第二步：列方程.}
\[
e^{0}=\alpha_0,\qquad
e^{\ii x}=(\ii x)^2\alpha_2+\ii x\alpha_1+\alpha_0,\qquad
e^{-\ii x}=(-\ii x)^2\alpha_2-\ii x\alpha_1+\alpha_0 .
\]
\textbf{第三步：解方程.} 由第一式 $\alpha_0=1$。后两式相减、相加：
\[
\alpha_1x=\frac{e^{\ii x}-e^{-\ii x}}{2\ii}=\sin x,\qquad
\alpha_2x^2=\frac{e^{\ii x}+e^{-\ii x}-2}{-2}=1-\cos x .
\]
\textbf{第四步：写出 $e^{Ax}$.} 代入待定式：
\[
e^{Ax}=\pmat{
1 & -2(1-\cos x)+\sin x & -5(1-\cos x)\\
0 & -(1-\cos x)-2\sin x+1 & -5\sin x\\
0 & \sin x & -(1-\cos x)+2\sin x+1}
\]
即
\[
e^{Ax}=\pmat{
1 & -2+2\cos x+\sin x & -5+5\cos x\\
0 & \cos x-2\sin x & -5\sin x\\
0 & \sin x & \cos x+2\sin x}.
\]
\textbf{自检.} (1) $x=0$ 得 $E$；(2) 求导在 $0$ 处等于 $A$；(3) $Y=e^{Ax}C$ 给出通解，$C$ 任意。

\subhdr{4.14 两条路线的对照}

\begin{center}
\begin{tabular}{p{0.18\linewidth}p{0.36\linewidth}p{0.38\linewidth}}
\hline
 & \thead{对角化} & \thead{Cayley--Hamilton}\\
\hline
前提 & 每个特征值 $g_j=m_j$，即有 $n$ 个线性无关特征向量 & 任意方阵；不需要可对角化\\
主要公式 & $e^{Ax}=Ve^{\Lambda x}V^{-1}$ & $e^{Ax}=\sum_{k=0}^{n-1}\alpha_k(\Ax)^k$\\
需要算 & 特征值、特征向量、$V^{-1}$ & 特征值；解 $n$ 元线性方程组（重根加导数条件）\\
优点 & 形式直观，列向量就是基本解 & 对重特征值也适用，不必求特征向量\\
易错点 & 特征向量不够仍强行求 $V^{-1}$；漏验 $AV=V\Lambda$ & 把 $\Ax$ 的特征值当成 $A$ 的；重根漏导数条件\\
\hline
\end{tabular}
\end{center}

两条路线给出的 $e^{Ax}$ 必须相同。若用两种方法各算一遍，$x=0$ 与求导自检可以互相印证；
这也是检查代数错误的有效办法。

\flowerrule

%=====================================================================
\clearpage
\hdr{§5}{基本矩阵（p39--p43）}
%=====================================================================

\subhdr{5.1 基本解组与基本矩阵}

前几节都假定 $A$ 是常数矩阵。课件 p40 先对更一般的线性系统给出定义：
\begin{equation}\label{eq:general}
Y'(x)-A(x)Y=0,
\end{equation}
其中 $A(x)$ 可以是 $x$ 的矩阵函数。

\callout{
\textbf{定义（p40）.} 若 $\{Y_1(x),\dots,Y_n(x)\}$ 是 \eqref{eq:general} 的 $n$ 个线性无关解，
则称它为一个\textbf{基本解组}，并称以它们为列的矩阵
\[
\Phi(x)=\bigl(Y_1(x),\dots,Y_n(x)\bigr)
\]
为 \eqref{eq:general} 的\textbf{基本矩阵}。此时通解为
$Y(x)=c_1Y_1(x)+\cdots+c_nY_n(x)=\Phi(x)C$，$C=(c_1,\dots,c_n)^T$ 任意。}

“线性无关”可以等价地写成 $\det\Phi(x)\ne0$（在 $A(x)$ 连续的区间上，Wronskian 要么恒为零、
要么处处非零）。写 $\Phi(x)C$ 时要注意：$C$ 是常向量，$\Phi$ 的列才是解；不要把 $C$ 放到左边。

\clearpage
\subhdr{5.2 常数矩阵情形：$e^{Ax}$ 与 $\Phi$ 的关系}

对常数 $A$，系统是
\begin{equation}\label{eq:H1}
Y'(x)-AY=0 .
\end{equation}
由 §1，通解可写成 $Y=e^{Ax}C$；由 5.1，通解又可写成 $Y=\Phi(x)C$。于是自然要问：
$e^{Ax}$ 与 $\Phi(x)$ 是什么关系（课件 p41 的 Q）。

课件 p42 的推导如下。考虑初值问题 $Y'=AY$、$Y(x_0)=Y_0$。因为
$Y(x_0)=\Phi(x_0)C=Y_0$ 且 $\det\Phi(x_0)\ne0$，所以 $C=\Phi^{-1}(x_0)Y_0$，从而
\[
Y(x)=\Phi(x)C=\Phi(x)\Phi^{-1}(x_0)Y(x_0)
=\bigl(\Phi(x)\Phi^{-1}(x_0)\bigr)Y_0 .
\]
另一方面，同一个初值问题的解在 §1 已写成 $Y=e^{A(x-x_0)}Y_0$。由初值问题解的唯一性
（两个解在 $x_0$ 处相同，且都满足同一线性方程），
\begin{equation}\label{eq:PhiExp}
\Phi(x)\Phi^{-1}(x_0)=e^{A(x-x_0)},\qquad
\Phi(x)\Phi^{-1}(0)=e^{Ax}.
\end{equation}
特别地，若 $\Phi(0)=E$，则 $e^{Ax}=\Phi(x)$。这说明：$e^{Ax}$ 是一个“在 $0$ 处取单位矩阵”的
基本矩阵；一般的基本矩阵 $\Phi$ 与它相差一个右乘常矩阵 $\Phi^{-1}(0)$。

\subhdr{5.3 Theorem 4.1：$e^{Ax}$ 的列构成基本解组}

\callout{
\textbf{Theorem 4.1（p43）.} 若 $A$ 是 $n\times n$ 常数矩阵，则 $e^{Ax}$ 的列构成
$Y'(x)=AY(x)$ 的一个基本解组；因此 $e^{Ax}$ 是该系统的基本矩阵。}

\textbf{证明.} 设 $e^{Ax}$ 的列为 $\Phi_1(x),\dots,\Phi_n(x)$。由 §1.4，$e^{Ax}$ 满足
$(e^{Ax})'=Ae^{Ax}$，逐列看就是 $\Phi_j'=A\Phi_j$，故每列都是解。再看 $x=0$：
\[
e^{A\cdot0}=E=(\varepsilon_1,\dots,\varepsilon_n),
\]
列向量在 $x=0$ 处是标准基，线性无关。若它们在某处线性相关，则由线性方程解空间的维数
（或 Wronskian 的性质）会在所有点线性相关，与 $x=0$ 处矛盾。因此 $n$ 列线性无关，
构成基本解组，$e^{Ax}$ 是基本矩阵。

\textbf{另一个等价说法.} $e^{Ax}$ 可逆，逆为 $e^{-Ax}$（§2 性质 (4)），所以
$\det e^{Ax}\ne0$；而它的列都是解，故自动是基本矩阵。这个说法更短，也顺便说明了
$\Phi(0)=E$ 时为什么 $e^{Ax}=\Phi(x)$。

\subhdr{5.4 教学小例：例题 5 的基本矩阵}

对 $A=\pmat{0&1\\-1&0}$，§4.8 已得
\[
e^{Ax}=\pmat{\cos x&\sin x\\-\sin x&\cos x}.
\]
它满足 $e^{A0}=E$，所以按 5.2，$\Phi(x)=e^{Ax}$ 就是一个基本矩阵；
两列 $(\cos x,-\sin x)^T$ 与 $(\sin x,\cos x)^T$ 线性无关（行列式为 $1$），
通解就是它们的线性组合。这个例子也说明：基本矩阵不唯一，但“在 $0$ 处等于 $E$”的那个唯一。

\flowerrule

%=====================================================================
\clearpage
\hdr{§6}{独立练习（课件 p44 作业 1--3）}
%=====================================================================

下面三道题取自课件 p44 的 Homework，题干与课件一致。它们要求\emph{同时}做两件事：
求 $e^{Ax}$，并求 $Y'=AY$ 的通解。请先独立完成，再翻到附录 A 核对；
正文只给题干与中性考点，不提前给出步骤、结果或谱分类。三题分别是 $2\times2$、$2\times2$、$3\times3$ 矩阵。

\practice{H1}{课件 p44 第 1 题}%
{设 $A=\pmat{-6&-4\\8&6}$。求 $e^{Ax}$，并求 $Y'=AY$ 的通解。}%
{2×2 常系数系统；矩阵指数；通解。}{ans:H1}

\practice{H2}{课件 p44 第 2 题}%
{设 $A=\pmat{-4&4\\-5&4}$。求 $e^{Ax}$，并求 $Y'=AY$ 的通解。}%
{2×2 常系数系统；矩阵指数；通解。}{ans:H2}

\practice{H3}{课件 p44 第 3 题}%
{设 $A=\pmat{-24&0&0\\4&-2&0\\1&0&2}$。求 $e^{Ax}$，并求 $Y'=AY$ 的通解。}%
{3×3 常系数系统；矩阵指数；通解。}{ans:H3}

\clearpage
\appendix
\hdr{附录 A}{练习详解}

\answerstart{H1}{课件 p44 第 1 题}%
{设 $A=\pmat{-6&-4\\8&6}$。求 $e^{Ax}$，并求 $Y'=AY$ 的通解。}

\textbf{第一步：特征值.}
\[
\det(A-\lambda E)=\begin{vmatrix}-6-\lambda&-4\\8&6-\lambda\end{vmatrix}
=(-6-\lambda)(6-\lambda)+32=\lambda^2-4=(\lambda-2)(\lambda+2).
\]
所以 $\lambda_1=2$、$\lambda_2=-2$，互不相同，$A$ 可对角化。

\textbf{第二步：特征向量.}
\[
\lambda=2:\quad (A-2E)V=\pmat{-8&-4\\8&4}V=0\ \Rightarrow\ 2x+y=0\ \Rightarrow\ V_1=\pmat{1\\-2};
\]
\[
\lambda=-2:\quad (A+2E)V=\pmat{-4&-4\\8&8}V=0\ \Rightarrow\ x+y=0\ \Rightarrow\ V_2=\pmat{1\\-1}.
\]
验算：$AV_1=(-6+8,\,8-12)^T=(2,-4)^T=2V_1$；
$AV_2=(-6+4,\,8-6)^T=(-2,2)^T=-2V_2$。

\textbf{第三步：对角化.}
\[
V=\pmat{1&1\\-2&-1},\qquad \Lambda=\pmat{2&0\\0&-2},\qquad
\det V=-1+2=1\ne0,
\]
\[
V^{-1}=\pmat{-1&-1\\2&1}
\quad(\text{验算 }VV^{-1}=E).
\]

\textbf{第四步：$e^{Ax}$.}
\[
e^{Ax}=V\pmat{e^{2x}&0\\0&e^{-2x}}V^{-1}
=\pmat{-e^{2x}+2e^{-2x}&-e^{2x}+e^{-2x}\\[2pt]2e^{2x}-2e^{-2x}&2e^{2x}-e^{-2x}}.
\]

\textbf{第五步：通解.}
\[
Y(x)=e^{Ax}C
=c_1\pmat{1\\-2}e^{2x}+c_2\pmat{1\\-1}e^{-2x},\qquad c_1,c_2\ \text{任意}.
\]
（这里 $(c_1,c_2)^T=V^{-1}C$；$C$ 任意等价于 $c_1,c_2$ 任意。）

\textbf{检查.} (1) $x=0$：$e^{A0}=\pmat{1&0\\0&1}$；(2) 求导后令 $x=0$：
\[
\left.\frac{\dd}{\dd x}e^{Ax}\right|_{0}
=\pmat{-2-4&-2-2\\4+4&4+2}
=\pmat{-6&-4\\8&6}=A;
\]
\answercontinue{H1}{H1 的自检 (3) 与 Cayley--Hamilton 交叉核对；$A=\pmat{-6&-4\\8&6}$，$e^{Ax}$ 见上一页。}
(3) 两列都满足 $Y'=AY$。
另外，用 Cayley--Hamilton 交叉核对：$n=2$ 时
$e^{Ax}=\alpha_1\Ax+\alpha_0E$，由 $e^{\pm2x}=\pm2x\alpha_1+\alpha_0$ 得
$\alpha_1x=\frac14(e^{2x}-e^{-2x})$、$\alpha_0=\frac12(e^{2x}+e^{-2x})$，
代回给出的矩阵与上式逐元素相同。

\answerstart{H2}{课件 p44 第 2 题}%
{设 $A=\pmat{-4&4\\-5&4}$。求 $e^{Ax}$，并求 $Y'=AY$ 的通解。}

\textbf{第一步：特征值.} 迹为 $0$，行列式为 $(-4)\cdot4-4\cdot(-5)=-16+20=4$，所以
\[
\det(A-\lambda E)=\lambda^2-(\operatorname{tr}A)\lambda+\det A=\lambda^2+4,
\]
特征值为 $\lambda=\pm2\ii$（对应 $\Ax$ 的特征值 $\pm2\ii x$），互不相同。

\textbf{第二步：用 Cayley--Hamilton 方法设待定式.} $n=2$，
\[
e^{Ax}=\alpha_1\Ax+\alpha_0E,\qquad r(\lambda)=\alpha_1\lambda+\alpha_0 .
\]
代入两个特征值 $\lambda_1=2\ii x$、$\lambda_2=-2\ii x$：
\[
e^{2\ii x}=2\ii x\alpha_1+\alpha_0,\qquad
e^{-2\ii x}=-2\ii x\alpha_1+\alpha_0 .
\]
两式相减、相加：
\[
\alpha_1x=\frac{e^{2\ii x}-e^{-2\ii x}}{4\ii}=\frac{\sin 2x}{2},\qquad
\alpha_0=\frac{e^{2\ii x}+e^{-2\ii x}}2=\cos 2x .
\]

\textbf{第三步：$e^{Ax}$.}
\[
e^{Ax}=\frac{\sin2x}{2}\pmat{-4&4\\-5&4}+\cos2x\pmat{1&0\\0&1}
=\pmat{\cos2x-2\sin2x&2\sin2x\\[2pt]-\frac52\sin2x&\cos2x+2\sin2x}.
\]

\textbf{第四步：通解.}
\[
Y(x)=c_1\pmat{\cos2x-2\sin2x\\[2pt]-\frac52\sin2x}
+c_2\pmat{2\sin2x\\[2pt]\cos2x+2\sin2x},\qquad c_1,c_2\ \text{任意}.
\]
若想用复特征向量写，也可取 $V=(2,2+\ii)^T$ 满足 $AV=2\ii V$，通解为
$\operatorname{Re}$ 与 $\operatorname{Im}$ 的实线性组合；与上式张成同一解空间。

\textbf{检查.} (1) $x=0$ 得 $E$；(2) 求导在 $0$ 处：
$(1,1)$ 为 $-4$，$(1,2)$ 为 $4$，$(2,1)$ 为 $-5$，$(2,2)$ 为 $4$，即 $A$；
(3) 两列都满足方程。也可用 $\operatorname{tr}A=0$、$\det A=4$ 复核特征值。

\answerstart{H3}{课件 p44 第 3 题}%
{设 $A=\pmat{-24&0&0\\4&-2&0\\1&0&2}$。求 $e^{Ax}$，并求 $Y'=AY$ 的通解。}

\textbf{第一步：特征值.} $A$ 是下三角矩阵，$\det(A-\lambda E)$ 等于对角元之积
$(-24-\lambda)(-2-\lambda)(2-\lambda)$，所以特征值为
$\mu_1=-24$、$\mu_2=-2$、$\mu_3=2$，互不相同。

\textbf{第二步：特征向量.} 分别解 $(A-\mu_jE)V=0$。
\begin{itemize}
  \item $\mu_1=-24$：$A+24E=\pmat{0&0&0\\4&22&0\\1&0&26}$，方程
        $4x+22y=0$、$x+26z=0$。取 $x=286$，得 $y=-52$、$z=-11$，即
        $V_1=(286,-52,-11)^T$（任何非零倍数都可以）。
  \item $\mu_2=-2$：$A+2E=\pmat{-22&0&0\\4&0&0\\1&0&4}$，方程给出 $x=0$、$z=0$，
        $y$ 任意，取 $V_2=(0,1,0)^T$。
  \item $\mu_3=2$：$A-2E=\pmat{-26&0&0\\4&-4&0\\1&0&0}$，方程给出 $x=0$、$y=0$，
        $z$ 任意，取 $V_3=(0,0,1)^T$。
\end{itemize}
验算：$AV_1=(-24\cdot286,\;4\cdot286-2\cdot(-52),\;286+2\cdot(-11))^T
=(-6864,1248,264)^T=-24V_1$，故 $(A+24E)V_1=0$；其余两个显然。

\textbf{第三步：对角化.}
\[
V=\pmat{286&0&0\\-52&1&0\\-11&0&1},\qquad
\Lambda=\pmat{-24&0&0\\0&-2&0\\0&0&2},\qquad \det V=286\ne0,
\]
\[
V^{-1}=\pmat{\frac1{286}&0&0\\[2pt]\frac2{11}&1&0\\[2pt]\frac1{26}&0&1}.
\]
（第二行来自 $52/286=2/11$，第三行来自 $11/286=1/26$；可乘 $VV^{-1}$ 核对。）

\textbf{第四步：$e^{Ax}$.} 由 $e^{Ax}=Ve^{\Lambda x}V^{-1}$，逐列组合：
\[
e^{Ax}=\pmat{
e^{-24x}&0&0\\[3pt]
\frac2{11}\left(e^{-2x}-e^{-24x}\right)&e^{-2x}&0\\[3pt]
\frac1{26}\left(e^{2x}-e^{-24x}\right)&0&e^{2x}}.
\]

\answercontinue{H3}{H3 的通解与检查；$A=\pmat{-24&0&0\\4&-2&0\\1&0&2}$，$e^{Ax}$ 见上一页。}
\textbf{第五步：通解.}
\[
Y(x)=c_1\pmat{286\\-52\\-11}e^{-24x}
+c_2\pmat{0\\1\\0}e^{-2x}
+c_3\pmat{0\\0\\1}e^{2x},\qquad c_1,c_2,c_3\ \text{任意}.
\]
把 $c_1$ 换成 $c_1/286$ 也可把第一个向量写成 $(1,-\frac2{11},-\frac1{26})^T$。

\textbf{检查.} (1) $x=0$ 得 $E$；(2) 求导在 $0$ 处：对角元为 $-24,-2,2$，
$(2,1)$ 为 $\frac2{11}(-2+24)=4$，$(3,1)$ 为 $\frac1{26}(2+24)=1$，即 $A$；
(3) 三列都满足 $Y'=AY$。由于三个特征值互不相同，本题也可用 §3 的对角化直接得到；
若用 §4 的 Cayley--Hamilton，会得到同一个 $e^{Ax}$，但需解一个 $3\times3$ 线性方程组。

\clearpage
\hdr{附录 B}{来源核对、覆盖与勘误}

\subhdr{B.1 核对方法}

本讲义依据唯一原资料 \texttt{inputs/\allowbreak 16\_17\_matrix\_exponential.pdf}（44 个物理页）。
逐页读取文字层，并用字形坐标重建矩阵的列位置（幻灯片的矩阵在 PDF 中是一个个独立的字形，
直接抽取会丢失行列结构）；确认 44 页均有文字层、没有纯图片页。所有公式、例题与作业另用
符号计算（SymPy）独立重算，脚本与结果见同目录的 \texttt{scripts/} 与 \texttt{verification.txt}。
覆盖到页的对应表见同目录 \texttt{source-map.json}；下表是摘要。

\subhdr{B.2 逐页覆盖摘要}

\begin{center}
\begin{tabular}{p{0.13\linewidth}p{0.45\linewidth}p{0.32\linewidth}}
\hline
\thead{课件页} & \thead{内容} & \thead{本讲义落点}\\
\hline
p1--p3 & 标题、目录、章节页 & 封面、§0.1、§1 开头\\
p4--p6 & 标量 IVP、级数定义、求导验证、初值平移、级数不实用 & §1.1--1.7\\
p7 & 矩阵指数的 5 条性质与特征向量性质 & §2.1--2.3\\
p8--p12 & 特征方程、$V,\Lambda$、$e^{Ax}=Ve^{\Lambda x}V^{-1}$、通解 & §3.1--3.4\\
p13--p15 & Example 1（3×3，三个不同实特征值） & §3.5（教学示范）\\
p16--p19 & CH 定理、Theorem 3.2、多项式降阶 & §4.1--4.3\\
p20 & Example 2（多项式降阶） & §4.4（教学示范）\\
p21--p23 & 解析函数、$e^{Ax}$ 的插值方程、Example 3 & §4.5--4.6（教学示范）\\
p24--p25 & Example 4（实特征值 4，$-2$） & §4.7（教学示范）\\
p26--p27 & Example 5（复特征值 $\pm\ii$） & §4.8（教学示范）\\
p28--p32 & 重特征值、导数条件、Theorem 3.3、Example 6、四步总结 & §4.9--4.11（教学示范）\\
p33--p35 & Example 7（三重特征值 $-2$） & §4.12（教学示范）\\
p36--p38 & Example 8（$0,\pm\ii$） & §4.13（教学示范）\\
p39--p43 & 基本矩阵、与 $e^{Ax}$ 的关系、Theorem 4.1 & §5.1--5.4\\
p44 & Homework 1--3 & §6 与附录 A\\
\hline
\end{tabular}
\end{center}

§6 的三道练习只给题干与中性考点（不含具体矩阵的谱分类），完整解答在附录 A；课件 8 个编号例题在正文中明示为教学示范，
均已给出完整过程，不作为未曝光练习使用。

\clearpage
\subhdr{B.3 勘误（按原页核对后判断）}

下表只列会影响公式或结果的地方；“原文写法”按 PDF 字形坐标还原，不靠记忆。
\begin{center}
\begin{tabular}{p{0.06\linewidth}p{0.28\linewidth}p{0.27\linewidth}p{0.29\linewidth}}
\hline
\thead{页} & \thead{原印刷内容} & \thead{应改为} & \thead{判断依据}\\
\hline
p4 & $x=x_0e^{atx}$ & $x=x_0e^{at}$ & 代回 $x'=ax$ 不成立；标量 IVP 的解唯一\\
p14 & $V^{-1}$ 的 $(1,2)$ 元为 $0$ & $-1$ & $VV^{-1}\ne E$；代数余子式；$x=0$ 自检\\
p25 & $\tilde c_1=3c_1+c_2$，$\tilde c_2=3c_1-c_2$ & $\tilde c_1=\frac{3c_1+c_2}{6}$，$\tilde c_2=\frac{3c_1-c_2}{6}$ & 矩阵等式含公共因子 $1/6$；$x=0$ 给 $(c_1,c_2)$\\
p27 & 第二解向量 $(-\sin x,\cos x)^T$ & $(\sin x,\cos x)^T$ & $e^{Ax}$ 的第二列；求导自检\\
p35 & $\tilde c_3=c_3-2c_2$ & $\tilde c_3=c_3-c_2$ & $x=0$ 的第三分量应为 $c_3$\\
p36 & $e^{Ax}$ 的 $(3,2)$ 元为 $-\alpha_1x$ & $+\alpha_1x$ & 按 $A^2,\Ax$ 逐项乘加；p38 最终结果\\
\hline
\end{tabular}
\end{center}

另外三处是记号/措辞问题，不影响结论，讲义已按正确形式表述：
\begin{itemize}
  \item p21 末句把系数写成 $\alpha_0,\dots,\alpha_n$；应为 $\alpha_0,\dots,\alpha_{n-1}$（共 $n$ 个）。
  \item p28 说重特征值的“前 $m-1$ 阶导数”为零；准确说法是
        $\Delta^{(j)}(\lambda_i)=0$ 对 $j=0,1,\dots,m-1$ 成立，从而
        $f^{(j)}(\lambda_i)=r^{(j)}(\lambda_i)$（$j=0,\dots,m-1$）。
  \item p32 Step 3 写作 $r^{m_1-1}(\lambda_1)$；应理解为 $r^{(m_1-1)}(\lambda_1)$（第 $m_1-1$ 阶导数）。
\end{itemize}

\subhdr{B.4 没有改动的内容}

p9 把重特征值按 $\lambda_1,\dots,\lambda_1,\lambda_{m_1+1},\dots,\lambda_n$ 排列，
p12 相应地出现 $c_2V_2e^{\lambda_1x}$，这不是笔误，而是“前 $m_1$ 个特征值都等于 $\lambda_1$”的
压缩写法；讲义 §3.1 已把同一意思展开说明，原式保留。p5 的定理措辞、p19 的
“$k\ge n$”等属于表达方式，不影响公式，未改。

\end{document}
